\section{深度 CNN：VD-CNN 字符级文本分类}

\subsection{动机：向视觉学习}

Conneau et al. (2017) 提出了 \textbf{VD-CNN}（Very Deep CNN），灵感来自视觉领域的深度网络。当时存在一个有趣的对比：

\begin{itemize}
    \item \textbf{视觉}：已经使用 30--100 层的深度 CNN（如 ResNet）
    \item \textbf{NLP}：LSTM 模型通常只有 2--4 层，最多 8 层
\end{itemize}

此外，视觉模型从\textbf{原始像素}开始，而 NLP 模型从\textbf{词}开始。VD-CNN 的思路是：像视觉一样，从\textbf{原始字符}开始，使用深度 CNN 进行文本分类。

\subsection{VD-CNN 架构}

\begin{figure}[H]
\centering
\includegraphics[width=0.95\textwidth]{slides-images/slide-030.jpg}
\caption{VD-CNN 架构：从字符级输入开始，通过多层卷积块和局部池化逐步构建高层表示\protect\footnotemark}
\end{figure}
\footnotetext{Slides 第31页。}

VD-CNN 的架构设计与视觉网络（如 VGGNet、ResNet）非常相似：

\begin{enumerate}
    \item \textbf{输入}：1024 个字符，每个字符用 16 维嵌入表示
    \item \textbf{第一组卷积块}：64 个 size-3 滤波器 $\times$ 2 层（带残差连接）
    \item \textbf{Local Pooling}（pool/2）：序列长度减半 $\rightarrow$ 512
    \item \textbf{第二组卷积块}：128 个 size-3 滤波器 $\times$ 2 层
    \item \textbf{Local Pooling}：$\rightarrow$ 256
    \item \textbf{第三组卷积块}：256 个 size-3 滤波器 $\times$ 2 层
    \item \textbf{Local Pooling}：$\rightarrow$ 128
    \item \textbf{第四组卷积块}：512 个 size-3 滤波器 $\times$ 2 层
    \item \textbf{$k$-Max Pooling}（$k=8$）：保留 8 个最大激活值
    \item \textbf{全连接层}：fc(4096, 2048) $\rightarrow$ fc(2048, 2048) $\rightarrow$ fc(2048, nClasses)
\end{enumerate}

\begin{knowledgebox}{VD-CNN 的视觉网络设计模式}
VD-CNN 借鉴了视觉深度网络的三个关键设计：
\begin{itemize}
    \item \textbf{逐步下采样}：局部池化在每一阶段将序列长度减半，同时特征维度翻倍
    \item \textbf{残差连接}：允许梯度直接通过跳跃连接传播，缓解深度网络的训练困难
    \item \textbf{顶部全连接层}：在所有卷积层之后使用多个全连接层做最终分类
\end{itemize}
这种``逐层缩小空间分辨率、增加特征维度''的金字塔结构是深度 CNN 的标志性设计。
\end{knowledgebox}

\subsection{感受野分析}

在第四组卷积块处，每个 trigram 滤波器的感受野已经覆盖了约 24 个字符（约 6 个词的长度）。这意味着最高层的每个特征已经能感知到较长范围的上下文。

\subsection{实验结果}

\begin{figure}[H]
\centering
\includegraphics[width=0.95\textwidth]{slides-images/slide-032.jpg}
\caption{VD-CNN 实验结果：29 层模型在多个大规模文本分类数据集上取得了最优或接近最优的结果\protect\footnotemark}
\end{figure}
\footnotetext{Slides 第33页。}

VD-CNN 在以下数据集上进行了评估：
\begin{itemize}
    \item 新闻分类（AG News、Sogou News）
    \item 知识图谱分类（DBpedia）
    \item 情感分析（Yelp、Amazon Reviews）
\end{itemize}

关键发现：
\begin{itemize}
    \item 29 层模型一致优于 9 层和 17 层模型——\textbf{更深确实更好}
    \item 在 Yelp 和部分 Amazon 数据集上超越了当时的最优方法
    \item 仅使用字符级输入，无需预训练词向量
\end{itemize}

\begin{importantbox}{字符级 CNN 的意义}
VD-CNN 证明了可以\textbf{完全从字符级别出发}，通过深度卷积网络达到甚至超越词级别模型的性能。这消除了对预训练词向量的依赖，也天然地处理了未登录词（OOV）问题。这种从``原始信号''出发的设计哲学后来也影响了 Transformer 中子词分割（subword tokenization）的思路。
\end{importantbox}

\subsection{本章小结}

VD-CNN 将视觉领域深度 CNN 的设计模式成功迁移到文本分类。通过字符级输入、深层卷积（最多 29 层）、残差连接和逐层池化，在多个大规模数据集上取得了竞争力强的结果，证明了深度和原始信号输入的价值。

%% ========== 树递归神经网络 ==========

